%! Characteristics of Rational Functions — Unit 5, Lesson 5.0 — Lesson Plan
\documentclass[10pt]{article}
\usepackage{algebra2-article}
\usepackage{algebra2-boxes}
\usepackage{graphicx}   % -article does NOT load graphicx; needed for thumbnails/screenshots
\graphicspath{{images/}}
\newcommand{\TallMath}[1]{$\displaystyle #1\rule[-1.4em]{0pt}{3.2em}$}
\newcommand{\CourseName}{Algebra 2: Shepherd}
\newcommand{\SchoolYear}{2026--2027}
\newcommand{\MeetingLength}{55 minutes}
\newcommand{\UnitNumberName}{Unit 5: Rational Functions \quad}
\newcommand{\LessonNumberName}{Lesson 5.0: Characteristics of Rational Functions}

\begin{document}
\begin{center}
    {\Large\bfseries \CourseName: \SchoolYear \\
    {\small\bfseries \UnitNumberName \LessonNumberName}}
\end{center}

\begin{tcolorbox}[colback=forestbg, colframe=forest, boxrule=0.9pt, arc=2mm,
    left=2mm, right=2mm, top=1.5mm, bottom=1.5mm]
    \textbf{Primary Objective:} Students can read the graph of a rational function --- a
    \emph{fraction} of polynomials, and the first family of the year whose graph can come in
    \emph{pieces} --- and name its characteristics: its \emph{domain restrictions} (every input making
    the denominator zero), the \emph{equations} of its \emph{vertical} and \emph{horizontal
    asymptotes}, and any \emph{hole} (removable discontinuity), deciding between a wall and a hole by
    whether the offending factor \emph{cancels}; together with its \emph{domain}, \emph{range},
    \emph{intercepts} and \emph{zeros}, the intervals where it is increasing/decreasing and
    positive/negative, and its \emph{end behavior} --- which for a rational function \emph{is} its
    horizontal asymptote. Students also explain how a rational graph differs from the polynomial graphs
    of Unit~4.
    \\[2pt]
    \textbf{Standards (2023 VA SOL):} \textbf{A2.F.2} --- investigate and analyze the characteristics
    of rational functions graphically and algebraically. This Lesson~0 \emph{introduces}
    \textbf{A2.F.2h} (determine the equations of any vertical and horizontal asymptotes of a function
    using a graph or equation) --- the row the \S3 spine marks new for Unit~5, together with
    \emph{holes} and \emph{domain restrictions}. It revisits \textbf{A2.F.2a} (domain, range, zeros,
    and intercepts, \emph{including graphs with discontinuities} --- the clause that first bites here),
    \textbf{A2.F.2c} (increasing/decreasing intervals, now written as unions), and \textbf{A2.F.2g}
    (end behavior). \textbf{A2.F.2b} (compare and contrast the characteristics of function families)
    frames the closing polynomial-vs.-rational contrast. \emph{Slant/oblique asymptotes are not in the
    2023 SOL and are not taught or assessed in this course.}
\end{tcolorbox}

\renewcommand{\arraystretch}{1.4}
\begin{skillbox}[Priority Ideas \& Skills]{goldbox}
\begin{tabularx}{\linewidth}{@{}X|X@{}}
\textbf{Priority Ideas \& Skills} & \textbf{Key Understandings (the ``why'')} \\[2pt]
\hline
\vspace{-6pt}
\begin{itemize}[leftmargin=1.1em, nosep, after=\vspace{2pt}]
  \item Find all \textbf{domain restrictions} by setting the \textbf{denominator} equal to zero.
  \item Decide, with the \textbf{cancel test}, whether a restriction is a \textbf{vertical asymptote} or a \textbf{hole}, and give the hole's coordinates.
  \item Write the \textbf{equation} of each vertical asymptote, and of the \textbf{horizontal asymptote} --- from a graph, or from the \textbf{degree comparison} $n<m$ / $n=m$ / $n>m$.
  \item State \textbf{end behavior} using the horizontal asymptote.
  \item Read \textbf{domain}, \textbf{range}, \textbf{intercepts}/\textbf{zeros}, and the \textbf{increasing/decreasing} and \textbf{positive/negative} intervals, as \textbf{unions}.
  \item \textbf{Compare and contrast} rational graphs with polynomial graphs.
\end{itemize}
&
\vspace{-6pt}
\begin{itemize}[leftmargin=1.1em, nosep, after=\vspace{2pt}]
  \item \emph{The numerator makes zeros; the denominator makes trouble.} Division by zero is the single source of every new feature today.
  \item Whether a break is a \emph{wall} or a \emph{single missing point} is decided by one question: does the factor cancel?
  \item An asymptote is a \textbf{line}, so its answer is an \textbf{equation} --- ``$x=1$,'' never ``$1$.''
  \item A \emph{vertical} asymptote is a restriction on \emph{inputs} (untouchable). A \emph{horizontal} asymptote only describes \emph{outputs} at the far ends --- so a graph \emph{may} cross it.
  \item Missing inputs are why intervals must break at each vertical asymptote; ``$(-\infty,\infty)$'' is now often wrong.
  \item Continuity was a free gift from polynomials. Losing it is what makes this a new family.
\end{itemize}
\end{tabularx}
\end{skillbox}

\begin{skillbox}[{Vocabulary, Concepts, \& Theorems}]{sky}
\renewcommand{\arraystretch}{1.3}
\begin{tabularx}{\linewidth}{@{}l X@{}}
\textbf{Rational function} & \TallMath{R(x)=\frac{P(x)}{Q(x)}} with $P,Q$ polynomials and $Q(x)\neq0$. \\
\textbf{Domain restriction} & An \emph{excluded value}: an input making $Q(x)=0$, thrown out because division by zero is undefined. \\
\textbf{Vertical asymptote} & A line $x=a$ the graph approaches but never touches or crosses; outputs grow without bound near it. Arises when the factor $(x-a)$ \emph{does not} cancel. \\
\textbf{Horizontal asymptote} & A line $y=b$ the graph approaches as $x\to\pm\infty$ --- the end-behavior statement. From the degrees: $n<m\Rightarrow y=0$; $n=m\Rightarrow y=$ ratio of leading coefficients; $n>m\Rightarrow$ none. \\
\textbf{Hole (removable discontinuity)} & A single missing point, arising when a restriction's factor \emph{cancels}; its $y$-value comes from the \emph{simplified} expression. \\
\textbf{Discontinuous graph} & A graph with breaks --- it cannot be traced without lifting the pencil. Rational is the first such family in this course. \\
\textbf{Degree comparison} & $n=\deg P$, $m=\deg Q$; the three cases above are the whole horizontal-asymptote rule. \\
\end{tabularx}
\end{skillbox}

\begin{fixedskillbox}[Activate Prior Knowledge \& Spiral Review]{forestbg}
    The Warm-Up rehearses the three skills this lesson stands on, one per new idea: \textbf{(1)}
    solving \emph{denominator} $=0$, including two factorable quadratics ($x^2-9$, $x^2+x-6$) --- the
    Unit~4 factoring toolkit, reused as the domain-restriction engine; \textbf{(2)} evaluating
    $y=\frac{1}{x-1}$ in \emph{two} tables --- closing in on $x=1$ (outputs $-10,-100,100,10$, which
    blow up \emph{and} flip sign) and running out to $x=1001$ (outputs $0.1,0.01,0.001$, which settle
    down) --- so students meet vertical and horizontal asymptote behavior numerically before either is
    named; and \textbf{(3)} factoring and cancelling $\frac{x^2-4}{x-2}$, then substituting $x=2$ into
    the \emph{original} to get $\frac00$. Debrief all three aloud, but do \emph{not} supply the words
    ``asymptote'' or ``hole'' --- let the Hook graph name them.
\end{fixedskillbox}

\begin{skillbox}[Hook]{forestbg}
    Show the graph of $f(x)=\dfrac{x+2}{x-1}$ (asymptotes drawn dashed; $(-2,0)$ and $(0,-2)$
    labeled). Open with the observation that reframes the whole unit: \textbf{``Every graph we have
    read this year --- lines, parabolas, polynomials --- was one unbroken curve. This one is in two
    pieces. Which input is missing, and why?''} ($x=1$ makes the denominator zero, so $f(1)$ does not
    exist --- there is no point to plot.) Then \textbf{``What is the curve doing as it runs off to the
    far left and far right?''} (Closing in on the line $y=1$ from below and from above, never landing
    on it.) Contrast the two behaviors deliberately: \textbf{``Why can the graph never cross the
    dashed vertical line, but the dashed horizontal line is only about the far ends?''} Then plant the
    third idea with the Warm-Up's leftover: \textbf{``Warm-Up item~3 also had a missing input, but
    that graph did not fly off anywhere. What happened instead?''} (It is the line $y=x+2$ with one
    point punched out.) Land the point: \emph{a rational function is a fraction, and the denominator
    runs the show --- it decides which inputs are thrown out, whether the break is a wall or a single
    missing point, and what the graph settles down to at the far ends.}
\end{skillbox}

\begin{skillbox}[Lesson]{forestbg}
\begin{multicols}{2}
\textbf{1. What makes a function rational --- and where it is undefined.} A \textbf{rational
function} is a quotient of polynomials $R(x)=\frac{P(x)}{Q(x)}$. Since division by zero is undefined,
every input with $Q(x)=0$ is a \textbf{domain restriction}. State the one procedure the whole unit
rests on: \emph{set the denominator equal to zero and solve}. Do $\frac{x+2}{x-1}$ ($x\neq1$) and
$\frac{5}{x^2-9}$ ($x\neq\pm3$). Flag immediately that a zero \emph{numerator} is fine --- it is an
$x$-intercept, not a restriction. Because inputs are missing, the graph can be
\textbf{discontinuous}, and a break appears in exactly two ways.

\textbf{2. Vertical asymptotes --- the walls.} A \textbf{vertical asymptote} is a line the graph
approaches but never meets; the Warm-Up's $-100 \to 100$ jump is what that looks like numerically.
Insist on the \emph{equation} form: \textbf{``$x=1$, not $1$''} --- A2.F.2h asks for an equation, so
mark bare numbers wrong from today. Ask \textbf{``why can a graph never cross it?''} (that input is
not in the domain --- nothing to plot), and use the two branches to motivate writing the domain as a
\textbf{union}.

\columnbreak

\textbf{3. Horizontal asymptotes \& end behavior --- the level.} A \textbf{horizontal asymptote} is
the line the graph approaches as $x\to\pm\infty$; for a rational function it \emph{is} the
end-behavior statement --- arms flatten toward a level instead of running to $\pm\infty$. Build the
\textbf{degree comparison} table: $n<m\Rightarrow y=0$ ($\frac{4}{x^2+1}$); $n=m\Rightarrow y=$ the
ratio of leading coefficients ($\frac{x+2}{x-1}\Rightarrow y=1$); $n>m\Rightarrow$ \emph{none}
($\frac{x^2}{x+1}$ grows without bound). Then kill the classic overgeneralization: a graph
\textbf{may} cross a horizontal asymptote --- check $\frac{2x}{x^2+1}$ at $x=0$ and it sits right on
$y=0$.

\textbf{4. Holes --- when a restriction cancels.} On $\frac{x^2-4}{x-2}=\frac{(x-2)(x+2)}{x-2}=x+2$,
$x\neq2$, the graph is the \emph{line} $y=x+2$ with one point punched out: a \textbf{hole}. Give the
\textbf{cancel test} (factor cancels $\Rightarrow$ hole; factor stays $\Rightarrow$ vertical
asymptote), find the hole's $y$-value from the \emph{simplified} form, and hammer the habit that
matters all unit: \emph{cancelling never restores the input} --- read restrictions off the
\textbf{original} denominator. Close with the both-at-once case
$\frac{x-3}{x^2-9}=\frac{1}{x+3}$, $x\neq3$: a hole at $x=3$, a wall at $x=-3$.

\textbf{5. The full feature list, and the contrast.} Re-read every characteristic of $f$ (asymptotes,
domain, range --- note the range excludes $1$ because the graph never \emph{reaches} it --- zeros from
the numerator, $y$-intercept from $f(0)$, decreasing on \emph{each} branch, positive/negative
intervals, end behavior). Finish with the \textbf{A2.F.2b} table: polynomial (continuous, domain all
reals, no asymptotes, arms to $\pm\infty$) vs.\ rational (possibly discontinuous, restricted domain,
asymptotes, arms flattening). Preview Lesson~5.1, \emph{simplifying rational expressions}.
\end{multicols}
\end{skillbox}

\begin{skillbox}[Active Monitoring]{redbox}
Circulate during the notes practice and Tier work. Watch for and correct:
\begin{itemize}[leftmargin=1.2em, nosep]
  \item setting the \emph{numerator} to zero to find a vertical asymptote --- zeros come from the top, restrictions from the bottom;
  \item answering with a bare number (``$1$'') instead of the \emph{equation} of the line (``$x=1$'') --- the A2.F.2h requirement;
  \item announcing a vertical asymptote at \emph{every} excluded value without running the cancel test --- the signature error of the lesson;
  \item reading restrictions off the \emph{simplified} expression, which erases the hole (the same mistake that produces extraneous solutions in 5.6);
  \item writing one interval across a vertical asymptote (``decreasing on $(-\infty,\infty)$'') instead of one per branch;
  \item believing a graph can never cross a \emph{horizontal} asymptote --- and, conversely, thinking every rational function must \emph{have} one.
\end{itemize}
Cold-call prompts: ``Which part of the fraction did you set equal to zero, and why?'' ``Does that
factor cancel --- so is it a wall or a hole?'' ``Is that the asymptote, or the \emph{equation} of the
asymptote?'' ``Can the graph ever sit on that line? Which line?''
\end{skillbox}

\begin{skillbox}[Group Work \& Differentiation]{redbox}
\textbf{Reading Rational Graphs} activity (tiered; mirrors the notes).
\begin{multicols}{3}
\textbf{Tier R --- Remediate.} From one pre-drawn graph of $g(x)=\frac{2}{x+2}$ (one vertical
asymptote, horizontal asymptote $y=0$): count the pieces, name the missing input, write both asymptote
equations, give the domain as a union and the $y$-intercept, explain why there is \emph{no}
$x$-intercept (constant numerator), and give the decreasing and positive intervals.

\columnbreak
\textbf{Tier A --- Approaching.} Complete a seven-row feature table for two functions --- and only one
of them has asymptotes: $p(x)=\frac{x-2}{x+1}$ (wall $x=-1$, level $y=1$) beside
$q(x)=\frac{x^2-1}{x-1}$, whose graph is a \emph{line with a hole} at $(1,2)$. Then \emph{justify}:
factor $q$, and explain why $x=1$ gives a hole rather than a wall --- and why $1$ is still excluded
after cancelling.

\columnbreak
\textbf{Tier E --- Extension.} Equation-only work: a four-row table (restrictions / vertical asymptotes
/ hole / horizontal asymptote) over $\frac{3}{x-5}$, $\frac{2x}{x+4}$, $\frac{x^2-25}{x-5}$, and
$\frac{x+1}{x^2+3x+2}$ (two restrictions, one asymptote, one hole), plus ``which one has \emph{no}
horizontal asymptote, and why?'' Then interpret an average-cost model
$A(n)=\frac{6n+250}{n}$: the level $y=6$ as the true per-shirt cost, why the average never reaches it
($A(n)=6+\frac{250}{n}$), and why $n=0$ is meaningless in context.
\end{multicols}
\end{skillbox}

\begin{skillbox}[Individual Work \& Assessment]{redbox}
\textbf{Exit Ticket} (independent, no notes): from one graphed rational function
$f(x)=\frac{x+3}{x+1}$, write the \emph{equations} of the vertical and horizontal asymptotes, give the
domain as a union, the $x$- and $y$-intercepts, and the two decreasing intervals; then \textbf{explain}
why the graph can never cross $x=-1$ \emph{and} why $y=1$ is missing from the range; then an
\textbf{SOL-style multiple-choice} item on $h(x)=\frac{x-4}{x^2-16}$ whose distractors each break
exactly one idea (ignore the cancelling / swap which factor cancels / read restrictions off the
simplified form). Collect and sort into ``got it / bare-number slip / interval-and-union slip /
hole-vs-asymptote slip'' to decide whether Lesson~5.1 opens with a re-teach --- the last pile matters
most, since 5.1 is built on exactly that distinction.
\end{skillbox}

\begin{skillbox}[Reinforcement \& Extension]{goldbox}
\textbf{Homework:} a four-row equation-only feature table (two rows require factoring first, and one
has two restrictions but only one asymptote); a full feature read off a two-asymptote graph
$m(x)=\frac{x-4}{x-2}$ (including range and end behavior); a feature read off a \emph{line with a
hole}, $n(x)=\frac{x^2-x-6}{x-3}$; and one justification item on why a graph may cross a horizontal
asymptote but never a vertical one. \textbf{Extension:} interpret a drug-concentration model
$C(t)=\frac{5t}{t^2+1}$ --- horizontal asymptote $y=0$ as the drug clearing the bloodstream, the peak
at $t=1$, the fact that the graph \emph{sits on} its horizontal asymptote at $t=0$, and the
deliberately unsettling case of a rational function with \emph{no} vertical asymptote at all (its
domain $t\ge0$ comes from context, not algebra). \textbf{Preview:} Lesson~5.1, \emph{Simplifying
Rational Expressions \& Domain Restrictions} --- the algebra that produces the holes we just read, with
restrictions always taken from the \emph{original} denominator.
\end{skillbox}
\end{document}
